\begin{abstract}
Detecting antisemitism in social-media text is harder than generic hate-speech classification: the criticism-vs-hate boundary on Israeli policy is contested under the IHRA Working Definition; antisemitic discourse relies on coded language, so identity terms alone are unreliable signals; and the data is imbalanced across identity sub-populations with very different base rates. On the GoldStandard2024 dataset~\citep{jikeli2024goldstandard} (11{,}311 IHRA-annotated tweets, 4.79:1 imbalance; test split 1{,}060 / 187 positives) we study model behaviour jointly along three connected axes --- accuracy on the positive class, calibration \emph{per sub-group} rather than only on average, and counterfactual stability under identity-token swaps. We propose \textbf{Counterfactual Calibration Invariance (CCI v2)}, a training-time loss penalising the Jensen-Shannon divergence between predictions on a sentence and on its identity-token-swapped counterfactual within a religiously-symmetric vocabulary, regulated by an explicit symmetry filter. A short structural lemma shows that scalar post-hoc Temperature Scaling cannot reduce the hard counterfactual flip rate at the decision threshold $1/2$; reducing it therefore requires either a training-time intervention or a non-monotone post-hoc method. \textbf{Key findings:} (i)~\textbf{Multicalibration vs counterfactual invariance is in direct tension}: applying multicalibration~\citep{hebertjohnson2018multicalibration} post-hoc to a strong baseline reduces the hard flip rate by 8\% but \emph{increases the L\textsubscript{1}-norm soft flip rate by 96\%} (from 0.032 to 0.063); to our knowledge the first quantitative report of this trade-off on a text classifier. (ii)~CCI v2 reduces the hard counterfactual flip rate (CFR) by 49\% on DeBERTa-v3-base (184M), 55\% on DeBERTa-v3-large (440M), and 73\% on RoBERTa-large (355M, different family) versus a strong focal+identity-mask baseline at iso-architecture, with positive-class F\textsubscript{1} preserved or improved at scale (0.72 best). The base comparisons are paired-bootstrap significant against four baselines (Davani CLP, Model C raw, Model C+TS, Model C+multicalibration); the within-architecture transfer is significant on DeBERTa-large and RoBERTa-large for CFR\textsubscript{hard} and CFR\textsubscript{soft, L\textsubscript{1}} under Holm-Bonferroni at $\alpha = 0.05$. (iii)~Under the Dixon~\citep{dixon2018measuring} sum-of-deviations FNED, CCI v2 sits within seed noise of every base-architecture baseline: $0.793$ vs.\ Davani $\lambda=0.5$ $0.786$ ($\Delta=+0.007$, $p_{\text{boot}}=0.86$, not significant) and vs.\ same-architecture Model~C $0.797$ ($\Delta=-0.004$, $p_{\text{boot}}=1.0$). The FNED differences at large scale are also within seed noise ($+5.6\%$ on DeBERTa-large, $+7.3\%$ on RoBERTa-large, neither Bonferroni-significant). (iv)~A side observation: under focal loss with F\textsubscript{1}-based early stopping we document a \emph{bimodal} ECE distribution across six seeds, traced to the interaction between the F\textsubscript{1}-peak epoch and the calibration trajectory. Integrated Gradients suggests that CCI~v2's CFR gains arise mainly from pair-level counterfactual constraints rather than large additional token-level salience suppression. Code, all 96 trained checkpoints, the 2{,}117-pair counterfactual test manifest, and the Holm-Bonferroni-corrected paired-bootstrap significance suite are publicly released. The prompted-LLM and cross-dataset transfer comparisons (Mistral-7B, Qwen-2.5-7B, Llama-3.1-8B, Llama-3.3-70B; HateXplain-Jewish and ToxiGen-Jewish) are part of a parallel evaluation track and are reported separately as they complete.
\end{abstract}
